\section{采样设计：on-policy、尾部条件采样与旧 rollout}
\label{sec:kl-sampling}

总定理从一开始区分 $\nu$ 与 $\rho$。本节把已有结论用于实际采样设计，所有积分目标和 log-ratio 保持不变。

\subsection{重要性校正与最优尾部常控制}
\begin{corollary}[直接重要性采样]\label{thm:kl-importance}
总定理取 $a=0$，便得到
\[
 \widehat I=\frac{\nu(A)}{\rho(A)}f_A,\qquad
 \E_\rho\widehat I=I_\nu(f),\qquad
 \Var_\rho(\widehat I)
 =\sum_v\frac{\nu(v)^2}{\rho(v)}\|f_v\|^2-\|I_\nu(f)\|^2,
\]
其中求和只作用于 $\rho(v)>0$ 的非零贡献。对反向 KL 的数值和梯度，分别取 $(\nu,f)=(p,r)$、$(p,rs)$；对正向目标，分别取 $(q,-r)$、$(q,-s)$。
\end{corollary}

尾部常控制的近似函数为头部 $a_v=f_v$、尾部 $a_v=b$。本小节比较所有可能的常控制，统一要求 $\rho(v)>0$ 对每个尾部动作成立；精确头部允许 $\rho(v)=0$。这个条件保证二次型对任意 $b$ 都有定义。定义
\[
 h_f=\sum_{v\in S}p(v)f_v,\qquad
 \mu_f=\E_{p_\tail}f_A,\qquad w_v=p(v)/\rho(v).
\]
总定理直接给出
\begin{equation}\label{eq:kl-is-generic}
 \widehat I_\rho(b)=h_f+\tau_pb+J_Aw_A(f_A-b),
\end{equation}
\begin{equation}\label{eq:kl-is-generic-var}
 \Var_\rho(\widehat I_\rho(b))
 =\sum_{v\notin S}\frac{p(v)^2}{\rho(v)}\|f_v-b\|^2
   -\tau_p^2\|\mu_f-b\|^2.
\end{equation}
为了与各估计量的标签对应，标量代入 $f=r,b=\beta$ 的式子记为
\begin{equation}\label{eq:kl-is-scalar}
 \widehat K_\rho(\beta)=H+\tau_p\beta+J_Aw_A(r_A-\beta),
\end{equation}
向量代入 $f=rs$ 的式子记为
\begin{equation}\label{eq:kl-is-vector}
 \widehat g_\rho(b)=h_g+\tau_pb+J_Aw_A(r_As_A-b).
\end{equation}
式~\eqref{eq:kl-is-generic-var} 的两次代入就是它们各自的方差；两种值域无需再证明两遍重要性采样。

\begin{proposition}[任意采样测度下的最优尾部常控制]
记
\[
 W=\sum_{v\notin S}\frac{p(v)^2}{\rho(v)},\qquad
 U_f=\sum_{v\notin S}\frac{p(v)^2}{\rho(v)}f_v,\qquad T_f=\tau_p\mu_f.
\]
若 $W>\tau_p^2$，则
\[
 b_\rho^*=\frac{U_f-\tau_pT_f}{W-\tau_p^2}.
\]
若 $W=\tau_p^2$，尾部控制逐样本抵消，方差与 $b$ 无关。
\end{proposition}
\begin{proof}
展开式~\eqref{eq:kl-is-generic-var}，与 $b$ 有关的项为
\[
 -2\langle U_f-\tau_pT_f,b\rangle+(W-\tau_p^2)\|b\|^2.
\]
当二次项系数严格为正时配方即得最优值。由 Cauchy--Schwarz，
\[
 \tau_p^2
 \le\left(\sum_{v\notin S}\frac{p(v)^2}{\rho(v)}\right)
      \left(\sum_{v\notin S}\rho(v)\right)\le W.
\]
等号要求 $\rho$ 全部质量在尾部，且与 $p_\tail$ 成比例，此时 $\rho=p_\tail$、$J_Aw_A=\tau_p$，式~\eqref{eq:kl-is-generic} 的控制项直接抵消。当 $\rho=p$ 时，$W=\tau_p$、$U_f=T_f$，最优控制退化为 $\mu_f$。
\end{proof}

\subsection{尾部条件采样与真正的分层方差减少}
若能够单独从 $p_\tail$ 采样 $B_1,\ldots,B_n$，则
\[
 \widehat I_\tail=h_f+\frac{\tau_p}{n}\sum_{i=1}^n f_{B_i},
 \qquad
 \Var(\widehat I_\tail)=\frac{\tau_p^2}{n}\Var_{p_\tail}(f).
\]
证明是在总定理中取 $\rho=p_\tail$，再独立平均。常控制 $\beta$ 或 $b$ 逐样本完全抵消，故此时 ResOPD 型的尾部常量控制没有额外作用。相同 $n$ 下，on-policy 头尾控制的方差为
\[
 \frac1n\{\tau_p\Var_{p_\tail}(f)+P\tau_p\|\mu_f-b\|^2\}.
\]
它与条件尾部采样的差为
\[
 \frac{P\tau_p}{n}\{\Var_{p_\tail}(f)+\|\mu_f-b\|^2\}\ge0.
\]
这是数学上的同样本数比较。实际采样器是否能取得尾部条件样本、是否需要额外 teacher 查询和拒绝采样，决定成本比较。

更一般地，若对每个分区单元独立抽样 $n_C$ 次，则
\[
 \widehat I_{\mathrm{strat}}
 =\sum_CP_C\frac1{n_C}\sum_{i=1}^{n_C}f_{B_{C,i}},\qquad
 \Var(\widehat I_{\mathrm{strat}})
 =\sum_C\frac{P_C^2}{n_C}\Var_{p_C}(f).
\]
这类分层积分直接消去了单元之间的抽样质量波动。给定连续预算 $\sum_Cn_C=N$、定义 $\sigma_C^2=\Var_{p_C}(f)$，Cauchy--Schwarz 给出
\[
 \sum_C\frac{P_C^2\sigma_C^2}{n_C}
 \ge\frac{(\sum_CP_C\sigma_C)^2}{N},
 \qquad n_C^*=\frac{NP_C\sigma_C}{\sum_DP_D\sigma_D}.
\]
零方差单元可以精确返回其常值；整数预算需另外取整。这里优化的是协方差迹，向量情形也使用相同 $\sigma_C$。

\subsection{截断与旧 rollout 的偏差由残差决定}
若仍以单位权重使用头尾控制，实际 $A\sim\rho$，总偏差定理给出
\begin{equation}\label{eq:kl-wrong-measure-bias}
 \E_\rho[\widehat I(b)]-I_p(f)
 =\sum_{v\notin S}(\rho(v)-p(v))(f_v-b).
\end{equation}
令 $\tau_\rho=\rho(S^c)$、$\mu_{f,\rho}=\E_\rho[f_A\mid J_A=1]$，条件协方差分解给出
\begin{equation}\label{eq:kl-wrong-measure-var}
 \Var_\rho(\widehat I(b))
 =\tau_\rho\Var_{\rho,\tail}(f_A)
 +(1-\tau_\rho)\tau_\rho\|\mu_{f,\rho}-b\|^2.
\end{equation}
因此更低的实际抽样方差与非零目标偏差可以同时出现。

旧 rollout 应使用实际生成概率 $\rho$ 校正；当前完整 softmax 只定义目标 $p$。温度变换保留支持时，若 $\rho$ 已知便可校正。top-$p$/top-$k$ 删除未知尾部动作时，总定理中的缺失支持项
\[
 -\sum_{v:\rho(v)=0}p(v)(f_v-a_v)
\]
给出无法由有限重要性权重恢复的贡献。只有已精确覆盖的动作残差为零时，支持缺失才不影响该目标。

若把正确权重 $w_v$ 截断为 $\bar w_v$，额外偏差为
\[
 \sum_{v:\rho(v)>0}\rho(v)(\bar w_v-w_v)(f_v-a_v).
\]
这也是同一个偏差定理的特例。通过近似减小未知动作上的残差，是缓解权重误差的方法；它不改变对支持和生成测度的要求。

\subsection{正向 Top-$J$：只替换积分权重与函数}

使用完全相同的 $p,q,S,\tau_p,\tau_q$，定义正向头部
\[
 H^{\mathrm F}=-\sum_{v\in S}q(v)r_v,\qquad
 B^{\mathrm F}=\tau_q\KL(q_{\tail}\Vert p_{\tail}).
\]
逐项将尾部 log-ratio 拆成质量与条件形状，得到
\[
 K^{\mathrm F}=H^{\mathrm F}-\tau_q\ell_{\tail}+B^{\mathrm F}.
\]
因此正向 Truncated、Tail-aware、Normalized 三个数值目标分别为
\[
 H^{\mathrm F},\qquad
 H^{\mathrm F}-\tau_q\ell_{\tail},\qquad
 \frac{H^{\mathrm F}}{Q}+\log\frac{P}{Q}.
\]
相对完整正向 $K^{\mathrm F}$ 的偏差分别为
\[
 \tau_q\ell_{\tail}-B^{\mathrm F},\qquad
 -B^{\mathrm F},\qquad
 \frac{H^{\mathrm F}}{Q}+\log\frac{P}{Q}-K^{\mathrm F},
\]
且固定 prefix 下方差均为零。这里 $B^{\mathrm F}$ 与前文 $B$ 分别对应不同方向的尾部条件 KL，上标始终保留。

对标量尾部控制 $\beta$，完整正向 KL 的无偏校正为
\[
 \widehat K^{\mathrm F}_\rho(\beta)
 =H^{\mathrm F}+\tau_q\beta
   +J_A\frac{q(A)}{\rho(A)}(-r_A-\beta).
\]
令 $\mu_r^{\mathrm F}=\E_{q_{\tail}}[-r_A]$，则其方差为
\[
 \Var_\rho(\widehat K^{\mathrm F}_\rho(\beta))
 =\sum_{v\notin S}\frac{q(v)^2}{\rho(v)}(-r_v-\beta)^2
   -\tau_q^2(\mu_r^{\mathrm F}-\beta)^2.
\]
证明是定理~\ref{thm:kl-importance} 在尾部对积分分布 $q$ 的直接应用：修正项均值是 $\tau_q(\mu_r^{\mathrm F}-\beta)$，恰好补齐尾部。

梯度方面，定义 $h_g^{\mathrm F}=-\sum_{v\in S}q(v)s_v$。因为 $q$ 固定，逐项求导给出
\begin{align*}
 \nabla H^{\mathrm F}&=h_g^{\mathrm F},\\
 \nabla(H^{\mathrm F}-\tau_q\ell_{\tail})
 &=h_g^{\mathrm F}-\tau_q\bar s_{\tail},\\
 \nabla\left(\frac{H^{\mathrm F}}{Q}+\log\frac P Q\right)
 &=\frac{h_g^{\mathrm F}}{Q}+\frac{d_P}{P}.
\end{align*}
三者的梯度偏差分别是上述向量减去 $g^{\mathrm F}$，梯度方差均为零。这里 $\bar s_{\tail}=\nabla\log\tau_p$ 是 student 尾部平均 score；它不能替换为未知的 teacher 尾部平均 score。特别地，正向 Tail-aware 梯度偏差为
\[
 \sum_{v\notin S}q(v)s_v-\tau_q\bar s_{\tail}.
\]
正向头部梯度与反向头部梯度的差异来自积分权重：$H^{\mathrm F}$ 的权重 $q$ 冻结，因此没有反向 $\nabla H$ 中的 $d_P$ 项。

正向无偏向量控制同样使用原来的向量符号 $b$：
\[
 \widehat g^{\mathrm F}_\rho(b)
 =h_g^{\mathrm F}+\tau_qb
   +J_A\frac{q(A)}{\rho(A)}(-s_A-b).
\]
令 $\mu_g^{\mathrm F}=\E_{q_{\tail}}[-s_A]$，有
\[
 \Var_\rho(\widehat g^{\mathrm F}_\rho(b))
 =\sum_{v\notin S}\frac{q(v)^2}{\rho(v)}\lVert-s_v-b\rVert^2
  -\tau_q^2\lVert\mu_g^{\mathrm F}-b\rVert^2.
\]
所有正向估计的支持覆盖要求作用于未知的 $q$ 尾部。这样，BPO 的正向 student--anchor center 与 OPD 的反向 student--teacher KL 都使用同一套采样与误差定义。
